
\documentclass[11pt]{article}

\usepackage[T1]{fontenc}
\usepackage[margin=0.83in]{geometry}
\usepackage{amsmath,amssymb}
\usepackage{booktabs}
\usepackage{parskip}
\usepackage{titlesec}
\titlespacing*{\section}{0pt}{1.2ex}{0.2ex}
\usepackage{setspace}
\setstretch{1.0}

\title{Lab 12}
\author{Chuan Wei Heng}
\date{\today}

\begin{document}
\maketitle
\vspace{-1.6em}

\section*{3.2 Exercises}

We compared the SciPy direct solver and LU factorization for six square
matrices. The timings (seconds) separate LU factorization from solving:

\begin{center}
\small
\begin{tabular}{rcccc}
\toprule
$N$ & Direct solve & LU factorization & LU solve & Break-even RHS \\
\midrule
100  & 0.0036 & 0.0005 & 0.0001 & 1 \\
500  & 0.0823 & 0.0555 & 0.0006 & 1 \\
1000 & 0.0177 & 0.0198 & 0.0032 & 2 \\
2000 & 0.1057 & 0.0572 & 0.0022 & 1 \\
4000 & 0.9368 & 0.4713 & 0.0053 & 1 \\
5000 & 1.8065 & 0.9918 & 0.0178 & 1 \\
\bottomrule
\end{tabular}
\end{center}

With the same matrix and $k$ right-hand sides, LU reuse is faster when
$k>t_{\mathrm{factor}}/(t_{\mathrm{direct}}-t_{\mathrm{LU\ solve}})$.
The measured break-even point was one right-hand side except for
$N=1000$, where it was two.

\section*{3.4 Exercises}

For $A\in\mathbb{R}^{100\times50}$, we compared the normal equations
$A^TAx=A^Tb$ with QR ($Rx=Q^Tb$). We chose $x_{\mathrm{true}}=\mathbf{1}$
and $b=Ax_{\mathrm{true}}$ to evaluate residuals and solution errors.

\begin{center}
\small
\begin{tabular}{llccc}
\toprule
Smallest $d$ & Method & Time (s) & Relative residual & Relative error \\
\midrule
$10^{-10}$ & Normal equations & 0.0228 & $7.65\times10^{-8}$ & $5.61$ \\
           & QR               & 0.0634 & $3.61\times10^{-16}$ & $5.56\times10^{-9}$ \\
$10^{-12}$ & Normal equations & 0.0011 & $2.82\times10^{-7}$ & $21.0$ \\
           & QR               & 0.0033 & $3.58\times10^{-16}$ & $2.98\times10^{-7}$ \\
\bottomrule
\end{tabular}
\end{center}

Reducing the smallest singular value increased $\kappa_2(A)$ from
$10^9$ to $10^{11}$. The normal equations produced large errors and
ill-conditioning warnings; QR was much more accurate because it avoids
forming $A^TA$, which squares the condition number.

\section*{Conclusion}

LU reuse can save time, while QR is more reliable for ill-conditioned
least-squares problems. The matrices are random, so measured results
can vary between runs. 

For details of the code, we turn the reader towards the Git repository by chuanweiheng-cell
entitled APPM4600. Thank you!

\end{document}
